% ===================================================================
%  CUADRO N.º 10 — El Mar. — El Puerto.
%  Traduction 2026 d'apres texto/io, controlee sur texto/fr et
%  texto/en. Consigne : voir texto/es/10-cuadro-01.tex.
% ===================================================================

%%K t10-apar-1 apar
\VUsaut{4.0mm}
\VUtitre{15.0pt}{60}{CUADRO N.º 10}
\VUsaut{3.0mm}
\VUpk{11.4pt}{40}{El Mar. --- El Puerto.}
\VUsaut{3.0mm}

%%K t10-01-1 p
1. --- En verano, mientras unos van a buscar la calma y el fresco a la
montaña, otros prefieren ir a la orilla del mar. Eso es lo que hicimos el
año pasado, porque nos gusta mucho el \VUgras{océano} \textsuperscript{(1)}.

%%K t10-02-1 p
2. --- Por mi parte, siento un gran placer al contemplar esa inmensa capa
de agua que se extiende hasta el \VUgras{horizonte} \textsuperscript{(2)} y al ver
partir para una larga \VUgras{travesía} los enormes \VUgras{barcos de vapor}
\textsuperscript{(3)} en los que embarcan los viajeros que van a América.

%%K t10-03-1 p
3. --- Por eso mi padre, que conoce mis gustos, elige siempre, para pasar
las vacaciones, una playa muy tranquila situada cerca de uno de nuestros
grandes \VUgras{puertos de guerra}.

%%K t10-03-2 p
Durante nuestra última estancia, la \VUgras{escuadra} había venido a echar
el ancla detrás del enorme \VUgras{dique} \textsuperscript{(4)} que cierra y protege
el \VUgras{puerto militar} \textsuperscript{(5)}, y pude admirar a mi gusto los
distintos \VUgras{buques de guerra}, desde el pequeño \VUgras{submarino}
\textsuperscript{(10)}, que se sumerge y desaparece bajo las aguas, hasta el enorme
\VUgras{acorazado} \textsuperscript{(9)}, que hasta ahora casi solo temía los ataques
del \VUgras{torpedero} \textsuperscript{(8)}.

%%K t10-tit-2 sub
\VUsaut{3.0mm}
\VUpk{10.2pt}{40}{Las embarcaciones pequeñas.}
\VUsaut{3.0mm}

%%K t10-04-1 p
4. --- Este ya sabía descubrir muy \VUgras{hábilmente} el punto débil de la
gruesa coraza metálica para clavar allí el peligroso \VUgras{torpedo}, cuya
explosión causa la pérdida del buque; pero era menos temible que el
submarino, porque los contratorpederos le daban caza fácilmente.

%%K t10-05-1 p
5. --- También pude ver los rápidos \VUgras{cruceros} \textsuperscript{(6)}
acorazados, casi tan invulnerables como el verdadero acorazado, varios
\VUgras{transportes} \textsuperscript{(12)}, tres o cuatro \VUgras{cañoneros}
\textsuperscript{(7)} y, por último, una \VUgras{fragata} \textsuperscript{(11)}.
\VUgras{El espectáculo de esa flota, cuando maniobra para levar anclas, es
de lo más interesante.}

%%K t10-06-1 p
6. --- ¡Qué diferencia de construcción entre esas enormes masas de acero y
los elegantes \VUgras{bergantines de tres palos} o los finos \VUgras{veleros}
\textsuperscript{(13)}, que se usan todavía en la marina mercante, y que yo veía
entrar en el puerto a toda vela cuando iba a pasear por el
\VUgras{espigón} \textsuperscript{(14)}! Es cierto que estos perdían mucho de sus
ventajas cuando el \VUgras{viento} era contrario. Tenían que recoger todas
las velas para entrar en la dársena, y hacía falta un \VUgras{remolcador}
\textsuperscript{(16)} que fuera a buscarlos para llevarlos al muelle como simples
\VUgras{gabarras} \textsuperscript{(17)}.

%%K t10-tit-3 sub
\VUsaut{3.0mm}
\VUpk{10.2pt}{40}{El puerto comercial.}
\VUsaut{3.0mm}

%%K t10-07-1 p
7. --- Lo interesante del puerto del que hablo es que no comprende solo un
puerto de guerra, creado artificialmente con obras gigantescas, sino
además un maravilloso \VUgras{puerto de comercio} situado en la
\VUgras{desembocadura} de un gran río.

%%K t10-08-1 p
8. --- La lengua de tierra que separa las dos dársenas está fortificada y
defendida por una serie de \VUgras{baterías} y \VUgras{baluartes} que se
adentran en el mar. Hace falta un permiso especial para pasear por las
\VUgras{murallas} \textsuperscript{(15)}. Yo lo había obtenido sin dificultad por
mediación de mi tío, que es ingeniero de los \VUgras{astilleros}
\textsuperscript{(18)}, y fui a visitar los barcos que se construyen bajo vastos
cobertizos.

%%K t10-09-1 p
9. --- Hasta asistí a la botadura de un acorazado y recuerdo el momento de
emoción que sintió toda la multitud cuando la enorme masa, abandonada a
sí misma, empezó a deslizarse por su grada y vino a flotar sobre el agua
del río.

%%K t10-10-1 p
10. --- Para ir a los astilleros se cruza el \VUgras{campo de maniobras}
\textsuperscript{(19)}, donde las tropas de la guarnición vienen cada día a hacer
el ejercicio, y se pasa por una \VUgras{pasarela} \textsuperscript{(20)} de hierro que
une la explanada con el \VUgras{arsenal} \textsuperscript{(21)}.

%%K t10-tit-4 sub
\VUsaut{3.0mm}
\VUpk{10.2pt}{40}{El Arsenal y los diques.}
\VUsaut{3.0mm}

%%K t10-11-1 p
11. --- Con mi tío visité ese gran establecimiento lleno de \VUgras{cañones}
\textsuperscript{(22)}, \VUgras{balas} \textsuperscript{(23)} y \VUgras{obuses}. Vi también la
\VUgras{fábrica metalúrgica} \textsuperscript{(24)} en la que se fabrican las
\VUgras{calderas} \textsuperscript{(25)} de los barcos de vapor.

%%K t10-12-1 p
12. --- Muy cerca están los \VUgras{diques} \textsuperscript{(26)} --- diques de
carena --- y los curiosísimos \VUgras{diques flotantes} \textsuperscript{(27)}, en
los que pude ver de cerca, en un barco en reparación, la \VUgras{hélice}
\textsuperscript{(28)}, la \VUgras{quilla} \textsuperscript{(29)} y el \VUgras{timón}
\textsuperscript{(30)} de una nave.

%%K t10-13-1 p
13. --- El puerto de comercio es tan interesante de estudiar como el
puerto de guerra, y pasé muchas horas callejeando por los muelles donde
están amarrados los \VUgras{vapores de ruedas} \textsuperscript{(31)} que remontan el
río, y mirando funcionar la \VUgras{grúa de arbolar} \textsuperscript{(32)}, que
levanta como plumas cargas enormes.

%%K t10-tit-5 sub
\VUsaut{4.0mm}
\VUpk{10.2pt}{40}{El Práctico.}
\VUsaut{3.0mm}

%%K t10-14-1 p
14. --- Admiraba los \VUgras{transatlánticos} \textsuperscript{(34)} que salían de la
rada con sus \VUgras{banderas} \textsuperscript{(35)} al viento, conducidos por un
hábil \VUgras{práctico} \textsuperscript{(36)} que sabía maravillosamente seguir el
\VUgras{canal} marcado por \VUgras{boyas} \textsuperscript{(40)} y \VUgras{balizas},
evitando las \VUgras{gabarras} \textsuperscript{(41)} que se dirigen con tanta
dificultad con su única vela cuadrada. Distinguía en la \VUgras{proa}
\textsuperscript{(37)} los \VUgras{camarotes} \textsuperscript{(39)} de segunda clase y, en la
\VUgras{popa} \textsuperscript{(38)}, los salones de los pasajeros de primera.

%%K t10-15-1 p
15. --- A veces asistía también a la carga de un \VUgras{buque de carga}
\textsuperscript{(42)}, en el que se amontonaban las mercancías del
\VUgras{almacén} \textsuperscript{(43)} y de la \VUgras{estación marítima}
\textsuperscript{(44)}, traídas por los numerosos trenes de la \VUgras{línea de los
muelles} \textsuperscript{(45)}.

%%K t10-tit-6 sub
\VUsaut{3.0mm}
\VUpk{10.2pt}{40}{Remar por gusto.}
\VUsaut{3.0mm}

%%K t10-16-1 p
16. --- A menudo remaba en la \VUgras{dársena a flote} \textsuperscript{(53)}, donde
se refugian las pequeñas \VUgras{embarcaciones} \textsuperscript{(46)}, y era una
alegría para mí manejar el \VUgras{remo} \textsuperscript{(47)}. Subía a la
\VUgras{cubierta} \textsuperscript{(48)} de una \VUgras{barca de vela} \textsuperscript{(49)},
trepaba por el \VUgras{palo} \textsuperscript{(51)} hasta las \VUgras{vergas}
\textsuperscript{(50)} con ayuda del \VUgras{cordaje} \textsuperscript{(52)}, y luego
reanudaba mi paseo en \VUgras{yola} \textsuperscript{(54)} entre las pesadas
\VUgras{barcas de pesca} \textsuperscript{(55)}, donde se secan las \VUgras{redes}
\textsuperscript{(56)}.

%%K t10-17-1 p
17. --- Por el \VUgras{muelle} \textsuperscript{(58)} desfilaban ante mí las
\VUgras{mercancías} \textsuperscript{(59)} más diversas, metidas en \VUgras{cajas}
\textsuperscript{(60)} o en \VUgras{fardos} \textsuperscript{(61)}. Formaban el
\VUgras{cargamento} \textsuperscript{(63)} de las barcazas, y potentes \VUgras{grúas}
\textsuperscript{(62)} las levantaban y las dejaban en \VUgras{camiones}
\textsuperscript{(64)} mientras los \VUgras{transportistas} hacían la declaración y
pagaban los derechos en la \VUgras{aduana} \textsuperscript{(65)}.

%%K t10-18-1 p
18. --- Por fin, cuando llegaba al \VUgras{puente giratorio} \textsuperscript{(66)} o
a la \VUgras{esclusa} \textsuperscript{(69)}, pasando por delante de la \VUgras{draga}
\textsuperscript{(68)} que limpia el \VUgras{canal} \textsuperscript{(67)}, desembarcaba en
el espigón, junto al \VUgras{semáforo} \textsuperscript{(70)}.

%%K t10-19-1 p
19. --- Es al otro lado de ese espigón donde atracan las \VUgras{chalupas}
\textsuperscript{(71)} que traen a tierra a los oficiales de la escuadra. Es un
espectáculo hermoso ver a los marineros levantar sus remos a compás
mientras la embarcación, llevada por la \VUgras{ola} \textsuperscript{(72)}, aborda
sin esfuerzo la escala del desembarcadero. Es aquí donde mejor se puede
observar la \VUgras{marea} y el movimiento del \VUgras{flujo} y del
\VUgras{reflujo} en la \VUgras{playa} \textsuperscript{(73)}.

%%K t10-tit-7 sub
\VUsaut{3.0mm}
\VUpk{10.2pt}{40}{El Casino de oficiales.}
\VUsaut{3.0mm}

%%K t10-20-1 p
20. --- Para volver a casa tomaba el \VUgras{tranvía eléctrico}
\textsuperscript{(74)}, cuya \VUgras{parada} \textsuperscript{(75)} está en el \VUgras{poste}
que hay delante del casino de los oficiales.

%%K t10-20-2 p
Desde el \VUgras{balcón} \textsuperscript{(76)} del casino, adornado con
\VUgras{anclas} \textsuperscript{(77)}, cañones y balas, veía a menudo una brillante
reunión de oficiales de marina: \VUgras{tenientes de navío} \textsuperscript{(78)}
con su espada, \VUgras{capitanes de artillería} \textsuperscript{(79)} y de
\VUgras{infantería} \textsuperscript{(80)} de marina con su \VUgras{sable}. Detrás de
ellos había un \VUgras{marinero} \textsuperscript{(81)} que les traía un
\VUgras{catalejo} cuando querían distinguir lo que pasaba a lo lejos.

%%K t10-21-1 p
21. --- Veía también a jóvenes \VUgras{alféreces de navío} \textsuperscript{(82)} que
recibían las órdenes de su \VUgras{comandante} \textsuperscript{(83)} o del
\VUgras{almirante} \textsuperscript{(84)}, mientras un \VUgras{cabo} \textsuperscript{(85)}
esperaba para llevarlas a bordo.

%%K t10-22-1 p
22. --- Desde ese balcón la vista es magnífica: se domina toda la playa
con sus \VUgras{tiendas} \textsuperscript{(86)} y sus \VUgras{casetas} \textsuperscript{(87)},
y se ve, a la hora del baño, a los \VUgras{bañistas} \textsuperscript{(88)} que
retozan alegremente delante del \VUgras{casino} \textsuperscript{(89)}.

%%K t10-23-1 p
23. --- En el extremo de la playa, la costa forma una pequeña
\VUgras{península} \textsuperscript{(91)} rocosa, donde viene a romper el
\VUgras{oleaje} \textsuperscript{(92)} y sobre la que está construido un \VUgras{faro}
\textsuperscript{(93)} que indica de noche el camino a los navegantes. Esa
península solo está unida a tierra por un \VUgras{istmo} \textsuperscript{(90)} muy
estrecho.

%%K t10-tit-8 sub
\VUsaut{3.0mm}
\VUpk{10.2pt}{40}{Vista general de la costa.}
\VUsaut{3.0mm}

%%K t10-24-1 p
24. --- El \VUgras{acantilado} \textsuperscript{(94)} se prolonga, recortado por el
mar, que dibuja en él caprichosamente una serie de \VUgras{bahías}
\textsuperscript{(95)} y \VUgras{promontorios} (cabos) que lo hacen de lo más
pintoresco.

%%K t10-24-2 p
En la \VUgras{meseta} \textsuperscript{(96)}, que domina el \VUgras{estrecho}
\textsuperscript{(98)}, hay un \VUgras{fuerte} \textsuperscript{(97)} muy importante y
algunos chalés.

%%K t10-25-1 p
25. --- Ese estrecho separa del continente una \VUgras{isla} \textsuperscript{(99)}
muy peligrosa, rodeada de \VUgras{escollos} \textsuperscript{(100)} y \VUgras{arrecifes},
donde barcas y hasta grandes buques encallan a veces. A pesar de la
prontitud de los socorros y de la rápida llegada de los \VUgras{botes
salvavidas}, esos \VUgras{naufragios} son terribles y, en las fuertes
\VUgras{tempestades} de equinoccio, varios barcos se han perdido con
hombres y mercancías.

%%K t10-25-2 p
A pesar de esa vecindad, \VUgras{el puerto} lo frecuentan mucho los marinos.
\VUsaut{6.0mm}
\VUfilet{22mm}
